% ch48_infty_yoneda.tex — Volume IV Chapter 12

\chapter{The $\infty$-Categorical Yoneda Lemma}

\begin{overviewbox}
The Yoneda lemma is the most read theorem of category theory. It has
appeared four times in this book: in Chapter~9 as a hom-set bijection;
in Chapter~15 as the representability principle; in Chapter~20 in
enriched form; in Chapter~28 unified via ends in the formal Yoneda
statement. This chapter delivers the fifth incarnation — the
$\infty$-categorical Yoneda lemma — which subsumes all of the above as
special cases.

The statement: for any (small) $\infty$-category $\mathcal{C}$, there is
a \term{Yoneda embedding}
\begin{equation*}
  y : \mathcal{C} \;\hookrightarrow\; \mathcal{P}(\mathcal{C}) := \mathrm{Fun}(\mathcal{C}^{\mathrm{op}}, \mathcal{S}),
\end{equation*}
where $\mathcal{S}$ is the $\infty$-category of spaces (anima). The
embedding is \emph{fully faithful}: it induces equivalences on mapping
spaces. The \term{$\infty$-Yoneda lemma} computes mapping spaces in
$\mathcal{P}(\mathcal{C})$ out of representables:
\begin{equation*}
  \mathrm{Map}_{\mathcal{P}(\mathcal{C})}(y(x), F) \;\simeq\; F(x), \quad \text{naturally in } x \in \mathcal{C} \text{ and } F \in \mathcal{P}(\mathcal{C}).
\end{equation*}
The structural consequence is that $\mathcal{P}(\mathcal{C})$ is the
\term{free cocompletion} of $\mathcal{C}$ — the universal $\infty$-category
containing $\mathcal{C}$ in which all small colimits exist.

This chapter is the milestone of Phase~2: the foundation is laid for
the rest of $\infty$-category theory. A reader who has completed
Chapter~48 can open Lurie's \emph{Higher Topos Theory} (especially HTT
§5.1, on presentable $\infty$-categories) or Riehl–Verity's
\emph{Elements of $\infty$-Category Theory} (their development of Yoneda
in Chapter~5) and read fluently.
\end{overviewbox}


% ═══════════════════════════════════════════════════════════════════════════
\section{Presheaves of Spaces}
\label{sec:presheaves-of-spaces}
% ═══════════════════════════════════════════════════════════════════════════

\subsection{Formalizing}

\begin{definition}[$\infty$-category of presheaves]
\label{def:presheaves}
For a small $\infty$-category $\mathcal{C}$, the
\term{$\infty$-category of presheaves of spaces} on $\mathcal{C}$ is
\begin{equation*}
  \mathcal{P}(\mathcal{C}) \;:=\; \mathrm{Fun}(\mathcal{C}^{\mathrm{op}}, \mathcal{S}),
\end{equation*}
the $\infty$-category of functors from $\mathcal{C}^{\mathrm{op}}$ to the
$\infty$-category of spaces $\mathcal{S}$.

Concretely, an object of $\mathcal{P}(\mathcal{C})$ is a functor of
quasi-categories $F : \mathcal{C}^{\mathrm{op}} \to \mathcal{S}$. A
morphism is a natural transformation, presented as a $1$-simplex of
$\mathcal{P}(\mathcal{C})$.
\end{definition}

\begin{howtosay}
$\mathcal{P}(\mathcal{C})$ \sayit{``the presheaves of spaces on
$\mathcal{C}$''}; sometimes $\widehat{\mathcal{C}}$ or $\mathrm{PSh}(\mathcal{C})$.
The dimensional shift is significant: presheaves of \emph{sets} (Chapter~9)
become presheaves of \emph{spaces} (Kan complexes / $\infty$-groupoids)
in the $\infty$-setting.
\end{howtosay}

\begin{remark}[Why presheaves of spaces, not of categories?]
\label{rem:why-spaces}
The natural choice for ``the values of a Yoneda embedding'' is the
$\infty$-category of \emph{$\infty$-groupoids} ($=$ spaces, $=
\mathcal{S}$), not the larger $\mathcal{C}\mathrm{at}_\infty$. Reason:
the representable functor $\mathrm{Map}_\mathcal{C}(-, x) :
\mathcal{C}^{\mathrm{op}} \to \mathcal{S}$ takes values in Kan complexes
— $\mathrm{Map}$ is always $\infty$-groupoidal (Theorem~44.\ref{thm:map-is-kan}).
We embed into $\mathcal{P}(\mathcal{C})$, not into a larger functor
category.
\end{remark}

\begin{theorem}[Properties of $\mathcal{P}(\mathcal{C})$]
\label{thm:presheaves-properties}
$\mathcal{P}(\mathcal{C})$ is:
\begin{enumerate}
  \item A presentable $\infty$-category (Chapter~52): cocomplete,
        accessible, and generated under small colimits by representables.
  \item Equipped with a Cartesian fibration to $\mathcal{C}\mathrm{at}_\infty^{\mathrm{op}}$
        (the universal fibration of Chapter~46), pulled back along the
        inclusion $\mathcal{S} \hookrightarrow \mathcal{C}\mathrm{at}_\infty$.
  \item Complete and cocomplete: $\infty$-(co)limits computed pointwise
        from $\mathcal{S}$.
\end{enumerate}
\end{theorem}

\begin{example_thm}[Presheaves on a $1$-category]
\label{ex:presheaves-1cat}
For $\mathcal{C} = \mathrm{N}(\mathcal{D})$ the nerve of an ordinary
category, $\mathcal{P}(\mathrm{N}(\mathcal{D}))$ is the
$\infty$-categorical refinement of the ordinary presheaf category
$\widehat{\mathcal{D}} = \mathbf{Set}^{\mathcal{D}^{\mathrm{op}}}$.
Specifically:
\begin{itemize}
  \item Objects of $\mathcal{P}(\mathrm{N}(\mathcal{D}))$ are functors
        $\mathcal{D}^{\mathrm{op}} \to \mathcal{S}$ — i.e., simplicial
        presheaves (presheaves of Kan complexes).
  \item Morphisms are natural transformations, $\infty$-categorically
        (homotopy-coherent).
\end{itemize}
$\mathrm{ho}(\mathcal{P}(\mathrm{N}(\mathcal{D})))$ recovers
$\mathrm{Ho}(\widehat{\mathcal{D}})$, the homotopy category of simplicial
presheaves.
\end{example_thm}


% ═══════════════════════════════════════════════════════════════════════════
\section{The $\infty$-Yoneda Embedding}
\label{sec:yoneda-embedding}
% ═══════════════════════════════════════════════════════════════════════════

\subsection{Formalizing}

\begin{definition}[$\infty$-Yoneda embedding]
\label{def:yoneda-embedding-qcat}
For a small $\infty$-category $\mathcal{C}$, the \term{$\infty$-Yoneda
embedding} is the functor
\begin{equation*}
  y : \mathcal{C} \;\longrightarrow\; \mathcal{P}(\mathcal{C})
\end{equation*}
characterized by: $y(x) = \mathrm{Map}_\mathcal{C}(-, x)$, the
representable presheaf at $x$.
\end{definition}

\begin{theorem}[Three constructions of $y$]
\label{thm:three-yoneda}
The Yoneda functor $y$ can be constructed in three equivalent ways:
\begin{enumerate}
  \item \textbf{Via straightening (Chapter~47).} The Cartesian fibration
        $\mathrm{ev}_1 : \mathcal{C}^{\Delta^1} \to \mathcal{C}$
        (Example~46.\ref{ex:target-eval}) restricts over each $x \in \mathcal{C}_0$
        to the slice $\mathcal{C}_{/x} \to *$. Straightening this fibration
        gives the functor $y : \mathcal{C} \to \mathcal{P}(\mathcal{C})$,
        $y(x) = \mathrm{Map}_\mathcal{C}(-, x)$.
  \item \textbf{Via the universal property.} $y$ is the unique
        colimit-preserving functor $\mathcal{P}(\mathcal{C}) \to
        \mathcal{P}(\mathcal{C})$ restricting to the identity on $\mathcal{C}$
        — i.e., $y$ sits inside the universal property of $\mathcal{P}(\mathcal{C})$
        as a free cocompletion (\S\ref{sec:free-cocompletion}).
  \item \textbf{Via cosimplicial spaces.} The cosimplicial space
        $\Delta^n \mapsto N(\mathcal{C}^{\mathrm{op}}_{/\bullet, [n]})$ — a
        specific cosimplicial simplicial set — Kan-extends along Yoneda
        to give $y$.
\end{enumerate}
All three are weakly equivalent in $\mathrm{Fun}(\mathcal{C}, \mathcal{P}(\mathcal{C}))$.
\end{theorem}

\begin{proof}[Proof sketch]
\emph{Construction (1).} The slice $\mathcal{C}_{/x}$ is a right fibration
(Example~46.\ref{ex:free-fib}). Its associated functor $\mathcal{C}^{\mathrm{op}}
\to \mathcal{S}$ under straightening (Chapter~47) is $\mathrm{Map}_\mathcal{C}(-, x)$.
Varying $x$ gives the functor $y : \mathcal{C} \to \mathcal{P}(\mathcal{C})$.

\emph{Construction (2)} is more conceptual: $\mathcal{P}(\mathcal{C})$
has a universal property as the free cocompletion (Theorem~\ref{thm:free-cocompletion}
below). The Yoneda embedding is forced by this universal property to
identify $\mathcal{C}$ with the generators.

\emph{Construction (3)} is the most concrete and computational, generalizing
Riehl–Verity's exposition. Each cosimplicial space gives a functor
$\mathbf{sSet}^{\mathcal{C}^{\mathrm{op}}} \to \mathbf{sSet}$, and the
specific cosimplicial $N(\mathcal{C}^{\mathrm{op}}_{/\bullet, [n]})$
recovers Yoneda.
\qed
\end{proof}


% ═══════════════════════════════════════════════════════════════════════════
\section{The $\infty$-Yoneda Lemma}
\label{sec:yoneda-lemma-statement}
% ═══════════════════════════════════════════════════════════════════════════

\subsection{Formalizing}

\begin{theorem}[$\infty$-Yoneda Lemma]
\label{thm:infty-yoneda}
For any small $\infty$-category $\mathcal{C}$ and any $F \in
\mathcal{P}(\mathcal{C})$, there is a weak equivalence of Kan complexes
\begin{equation*}
  \mathrm{Map}_{\mathcal{P}(\mathcal{C})}(y(x), F) \;\simeq\; F(x),
\end{equation*}
natural in $x \in \mathcal{C}$ and $F \in \mathcal{P}(\mathcal{C})$.
\end{theorem}

\begin{proof}[Proof sketch]
The strategy mirrors the $1$-categorical Yoneda proof (Chapter~9), but
all sets become spaces.

\emph{Construct the comparison map.} Given a natural transformation
$\eta : y(x) \to F$, take $\eta_x(\mathrm{id}_x) \in F(x)$ — the image of
the universal element. This defines $\Phi : \mathrm{Map}_{\mathcal{P}(\mathcal{C})}(y(x), F)
\to F(x)$.

\emph{Construct the inverse.} Given $\xi \in F(x)$, define
$\widetilde \xi : y(x) \to F$ pointwise: $\widetilde \xi_z : \mathrm{Map}_\mathcal{C}(z, x)
\to F(z)$ sends $f$ to $F(f)(\xi)$. Coherence of $\widetilde \xi$ as a
natural transformation requires $\infty$-categorical functoriality of $F$,
which is built in.

\emph{Verify mutual inverse.} The two compositions are: $\xi \mapsto
\widetilde \xi \mapsto \widetilde \xi_x(\mathrm{id}_x) = F(\mathrm{id}_x)(\xi)
\simeq \xi$ (functoriality). The other composition $\eta \mapsto \eta_x(\mathrm{id}_x)
\mapsto \widetilde{\eta_x(\mathrm{id}_x)}$ recovers $\eta$ by naturality.
The required coherences (at $2$-simplex level and above) are checked
combinatorially.
\qed
\end{proof}

\begin{corollary}[Yoneda for representables]
\label{cor:yoneda-representables}
For $x, y \in \mathcal{C}$:
\begin{equation*}
  \mathrm{Map}_{\mathcal{P}(\mathcal{C})}(y(x), y(y)) \;\simeq\; \mathrm{Map}_\mathcal{C}(x, y).
\end{equation*}
\end{corollary}

\begin{proof}
Apply $\infty$-Yoneda to $F = y(y)$:
\begin{equation*}
  \mathrm{Map}_{\mathcal{P}(\mathcal{C})}(y(x), y(y)) \simeq y(y)(x) = \mathrm{Map}_\mathcal{C}(x, y). \qed
\end{equation*}
\end{proof}

\begin{theorem}[Yoneda embedding is fully faithful]
\label{thm:yoneda-ff}
The functor $y : \mathcal{C} \to \mathcal{P}(\mathcal{C})$ is fully
faithful: $\mathrm{Map}_\mathcal{C}(x, y) \xrightarrow{\sim}
\mathrm{Map}_{\mathcal{P}(\mathcal{C})}(y(x), y(y))$ is a weak homotopy
equivalence of Kan complexes for all $x, y$.
\end{theorem}

\begin{proof}
By Corollary~\ref{cor:yoneda-representables} the comparison map is a
weak equivalence; this is exactly full faithfulness in the
$\infty$-categorical sense (Theorem~44.\ref{thm:whitehead-qcat}).
\qed
\end{proof}


% ═══════════════════════════════════════════════════════════════════════════
\section{Free Cocompletion}
\label{sec:free-cocompletion}
% ═══════════════════════════════════════════════════════════════════════════

\subsection{Formalizing}

\begin{theorem}[Universal property of $\mathcal{P}(\mathcal{C})$]
\label{thm:free-cocompletion}
For every cocomplete $\infty$-category $\mathcal{D}$, the restriction
functor
\begin{equation*}
  \mathrm{Fun}^{\mathrm{cc}}(\mathcal{P}(\mathcal{C}), \mathcal{D}) \;\xrightarrow{\;y^*\;}\; \mathrm{Fun}(\mathcal{C}, \mathcal{D})
\end{equation*}
(from colimit-preserving functors $\mathcal{P}(\mathcal{C}) \to \mathcal{D}$
to functors $\mathcal{C} \to \mathcal{D}$, via precomposition with the
Yoneda embedding) is an equivalence.

In words: $\mathcal{P}(\mathcal{C})$ is the \term{free cocompletion} of
$\mathcal{C}$. Every functor $F : \mathcal{C} \to \mathcal{D}$ to a
cocomplete $\infty$-category extends uniquely (up to coherent homotopy)
to a colimit-preserving functor $\bar F : \mathcal{P}(\mathcal{C}) \to
\mathcal{D}$.
\end{theorem}

\begin{proof}[Proof sketch]
\emph{Construction of $\bar F$.} Every $P \in \mathcal{P}(\mathcal{C})$
is a colimit of representables (\term{density theorem}, below). Set
$\bar F(P) := \mathrm{colim}\,(F \circ p)$ where $p : \mathcal{C}_{/P} \to \mathcal{C}$
is the canonical diagram of representables.

\emph{Uniqueness of $\bar F$.} A colimit-preserving functor is determined
by its restriction to a generating set; representables generate
$\mathcal{P}(\mathcal{C})$ under colimits.

\emph{Equivalence.} The equivalence is established by checking
fully-faithfulness and essential surjectivity of $y^*$, both of which
follow from the density theorem.
\qed
\end{proof}

\begin{theorem}[Density theorem]
\label{thm:density}
Every $P \in \mathcal{P}(\mathcal{C})$ is canonically a colimit of
representables: there is an equivalence
\begin{equation*}
  P \;\simeq\; \mathrm{colim}_{(x, \xi) \in \mathcal{C}_{/P}}\, y(x),
\end{equation*}
where $\mathcal{C}_{/P}$ is the slice of $\mathcal{C}$ over $P$
(via $y$), and the colimit is computed in $\mathcal{P}(\mathcal{C})$.
\end{theorem}

\begin{remark}[Foreshadowing Chapter~49]
The density theorem is the cornerstone of $\infty$-categorical limit
and colimit theory. It says: in $\mathcal{P}(\mathcal{C})$, representables
form a \emph{dense generator}. Constructing functors out of
$\mathcal{P}(\mathcal{C})$ that preserve colimits is the same as
specifying their values on representables — a powerful technique used
throughout Chapter~49 to construct $\infty$-(co)limits in general
$\infty$-categories.
\end{remark}


% ═══════════════════════════════════════════════════════════════════════════
\section{Comparison with the Earlier Yonedas}
\label{sec:yoneda-comparison}
% ═══════════════════════════════════════════════════════════════════════════

\subsection{Reorganizing}

\begin{remark}[The Yoneda quartet, revisited]
\label{rem:yoneda-quartet}
The $\infty$-Yoneda lemma (Theorem~\ref{thm:infty-yoneda}) is the
common ancestor of the four earlier appearances:
\begin{itemize}
  \item \textbf{Chapter~9 (ordinary Yoneda).} For $\mathcal{C}$ an
        ordinary category, $\mathcal{P}(\mathrm{N}(\mathcal{C}))$ contains
        $\widehat{\mathcal{C}} = \mathbf{Set}^{\mathcal{C}^{\mathrm{op}}}$
        as the full subcategory of $\pi_0$-valued presheaves (presheaves
        of \emph{discrete} spaces). Restricting $\infty$-Yoneda to this
        subcategory recovers $\mathrm{Hom}(y(x), F) \cong F(x)$.
  \item \textbf{Chapter~15 (representability).} The $\infty$-statement
        ``every presheaf of spaces is a colimit of representables''
        (Density theorem) is the higher version of Chapter~15's
        ``representability principle.''
  \item \textbf{Chapter~20 (enriched Yoneda).} For $\mathcal{V}$-enriched
        categories, the $\mathcal{V}$-Yoneda lemma is a special case of
        $\infty$-Yoneda by setting $\mathcal{C}$ to the $\infty$-category
        underlying $\mathcal{V}\mathrm{Cat}$.
  \item \textbf{Chapter~28 (formal Yoneda via ends).} The end formula
        $[\mathcal{C}, \mathcal{V}](y(-), F) \cong F$ in coend form
        generalizes to $\infty$-end form
        $\mathrm{Map}_{\mathcal{P}(\mathcal{C})}(y(-), F) \simeq F$
        — the precise $\infty$-statement we have just proved.
\end{itemize}
\end{remark}

\begin{remark}[Phase-2 milestone]
\label{rem:phase2-milestone}
With $\infty$-Yoneda in hand, the reader has the central organizing
theorem of $\infty$-category theory. Chapter~49 will use it to develop
$\infty$-(co)limits in general $\infty$-categories. Chapter~50 will
develop $\infty$-adjunctions, building on the representability that
$\infty$-Yoneda provides. The literature (Lurie HTT, Riehl–Verity, Land)
is now accessible.
\end{remark}


% ═══════════════════════════════════════════════════════════════════════════
\section{Exercises}
\label{sec:yoneda-exercises}
% ═══════════════════════════════════════════════════════════════════════════

\begin{exercisesbox}
\begin{qa}
\qaitem{What is $\mathcal{P}(\mathcal{C})$?}{$\mathrm{Fun}(\mathcal{C}^{\mathrm{op}}, \mathcal{S})$ — the $\infty$-category of presheaves of spaces on $\mathcal{C}$.}
\qaitem{Why values in $\mathcal{S}$ and not $\mathcal{C}\mathrm{at}_\infty$?}{Because $\mathrm{Map}_\mathcal{C}(-, x)$ takes values in Kan complexes (Theorem~44.2). For the Yoneda embedding to land in the natural target, presheaves of spaces is the right notion.}
\qaitem{State the Yoneda embedding.}{$y : \mathcal{C} \to \mathcal{P}(\mathcal{C})$, $y(x) = \mathrm{Map}_\mathcal{C}(-, x)$.}
\qaitem{Three constructions of $y$?}{Via straightening of the slice fibration; via universal property (free cocompletion); via cosimplicial-space Kan extension.}
\qaitem{State the $\infty$-Yoneda lemma.}{$\mathrm{Map}_{\mathcal{P}(\mathcal{C})}(y(x), F) \simeq F(x)$, natural in $x$ and $F$.}
\qaitem{What is the comparison map in the Yoneda proof?}{$\eta \mapsto \eta_x(\mathrm{id}_x)$ — evaluation at the universal element.}
\qaitem{What is the inverse map?}{$\xi \mapsto \widetilde \xi$ where $\widetilde \xi_z(f) = F(f)(\xi)$ — transporting $\xi$ along $F$'s contravariant action.}
\qaitem{State the corollary for representables.}{$\mathrm{Map}_{\mathcal{P}(\mathcal{C})}(y(x), y(y)) \simeq \mathrm{Map}_\mathcal{C}(x, y)$.}
\qaitem{Why does the corollary imply full faithfulness of $y$?}{Full faithfulness of an $\infty$-functor $=$ weak equivalence on each mapping space (Whitehead, Theorem~44.10). The corollary is exactly this for $y$.}
\qaitem{State the free cocompletion property.}{$\mathrm{Fun}^{\mathrm{cc}}(\mathcal{P}(\mathcal{C}), \mathcal{D}) \xrightarrow{y^*} \mathrm{Fun}(\mathcal{C}, \mathcal{D})$ is an equivalence — every $F : \mathcal{C} \to \mathcal{D}$ to cocomplete $\mathcal{D}$ extends uniquely to a colimit-preserving $\bar F$.}
\qaitem{State the density theorem.}{Every $P \in \mathcal{P}(\mathcal{C})$ is canonically the colimit of representables: $P \simeq \mathrm{colim}_{(x, \xi)} y(x)$ over $(x, \xi) \in \mathcal{C}_{/P}$.}
\qaitem{How does Chapter~9's Yoneda fit?}{For $\mathcal{C}$ a $1$-category, $\mathcal{P}(\mathrm{N}(\mathcal{C}))$ contains $\widehat{\mathcal{C}}$ as the $\pi_0$-valued subcategory; restricting $\infty$-Yoneda there recovers Chapter~9.}
\qaitem{How does Chapter~28's formal Yoneda fit?}{The end formula $\mathrm{Nat}(y(-), F) \cong F$ becomes $\mathrm{Map}_{\mathcal{P}(\mathcal{C})}(y(-), F) \simeq F$ in $\infty$-categorical form.}
\qaitem{What is the categorical content of ``free cocompletion''?}{$\mathcal{P}(\mathcal{C})$ is universal among cocomplete $\infty$-categories containing $\mathcal{C}$: any other such cocomplete embedding factors through it via a colimit-preserving functor.}
\qaitem{Why does the density theorem hold?}{Because every presheaf $P : \mathcal{C}^{\mathrm{op}} \to \mathcal{S}$ is determined by its values at all objects of $\mathcal{C}$ via Yoneda; the canonical colimit over $\mathcal{C}_{/P}$ reconstructs $P$.}
\qaitem{What is the slice $\mathcal{C}_{/P}$ when $P \in \mathcal{P}(\mathcal{C})$?}{The $\infty$-category of pairs $(x, \xi)$ where $x \in \mathcal{C}$ and $\xi \in P(x)$, with morphisms in $\mathcal{C}$ compatible with $P$'s action on $\xi$.}
\qaitem{What's the relationship to Cartesian fibrations?}{$\mathcal{C}_{/P} \to \mathcal{C}$ is a right fibration, classified under straightening by the presheaf $P$ itself.}
\qaitem{Why is $\mathcal{P}(\mathcal{C})$ presentable?}{It is cocomplete (pointwise from $\mathcal{S}$), accessible (small generated under filtered colimits), and locally small. Chapter~52 develops presentability.}
\qaitem{What property does Yoneda preserve?}{$y$ preserves all limits that exist in $\mathcal{C}$ (since representables preserve limits). It does \emph{not} in general preserve colimits — that's exactly the failure that makes $\mathcal{P}(\mathcal{C})$ the \emph{free} cocompletion.}
\qaitem{When does Yoneda preserve colimits?}{Never in general — that's the point of free cocompletion. It would only happen if $\mathcal{C}$ already had all colimits and they were preserved by the embedding, which is essentially never.}
\qaitem{What does ``representable presheaf'' mean?}{$y(x) = \mathrm{Map}_\mathcal{C}(-, x)$ for some $x \in \mathcal{C}$. A presheaf is representable iff it is in the essential image of $y$.}
\qaitem{What's the standard reading of the Yoneda philosophy?}{An object is determined by its incoming morphisms (its representable). Two objects are equivalent iff their representables are. This generalizes from $\mathbf{Cat}$ to $\mathcal{C}\mathrm{at}_\infty$ without change.}
\qaitem{Why is this chapter the milestone?}{Because with $\infty$-Yoneda in hand, the reader can use the foundation in any $\infty$-categorical reference (Lurie's HTT especially), which assumes Yoneda freely from the outset. The bridge to literature is now fully built.}
\end{qa}
\end{exercisesbox}


% ═══════════════════════════════════════════════════════════════════════════
\section*{Chapter Summary}
\addcontentsline{toc}{section}{Chapter Summary}
% ═══════════════════════════════════════════════════════════════════════════

\begin{summarybox}
\textbf{Presheaves of spaces.}\quad $\mathcal{P}(\mathcal{C}) =
\mathrm{Fun}(\mathcal{C}^{\mathrm{op}}, \mathcal{S})$. Cocomplete and
presentable; complete; the $\infty$-categorical refinement of ordinary
presheaf categories.

\textbf{Yoneda embedding.}\quad $y : \mathcal{C} \hookrightarrow
\mathcal{P}(\mathcal{C})$, $y(x) = \mathrm{Map}_\mathcal{C}(-, x)$.
Constructed via straightening of slice fibrations.

\textbf{$\infty$-Yoneda lemma.}\quad For all $x \in \mathcal{C}$, $F \in
\mathcal{P}(\mathcal{C})$:
$\mathrm{Map}_{\mathcal{P}(\mathcal{C})}(y(x), F) \simeq F(x)$,
natural in both. Corollary: $y$ is fully faithful.

\textbf{Free cocompletion.}\quad $\mathcal{P}(\mathcal{C})$ is universal:
$\mathrm{Fun}^{\mathrm{cc}}(\mathcal{P}(\mathcal{C}), \mathcal{D}) \simeq
\mathrm{Fun}(\mathcal{C}, \mathcal{D})$ for cocomplete $\mathcal{D}$.

\textbf{Density theorem.}\quad Every $P \in \mathcal{P}(\mathcal{C})$ is
canonically a colimit of representables.

\textbf{Bridge to literature.}\quad With $\infty$-Yoneda established,
Lurie HTT, Riehl–Verity, Cisinski, and Kerodon are now accessible. The
remaining $\infty$-categorical theorems of Volume~IV (limits,
adjunctions, monads, Kan extensions, presentability, stability, operads)
all derive from this foundation.
\end{summarybox}


% ═══════════════════════════════════════════════════════════════════════════
\section*{Formal Restatement}
\addcontentsline{toc}{section}{Formal Restatement}
% ═══════════════════════════════════════════════════════════════════════════

\begin{formalrestatement}
\textbf{Presheaves of spaces.}\quad $\mathcal{P}(\mathcal{C}) :=
\mathrm{Fun}(\mathcal{C}^{\mathrm{op}}, \mathcal{S})$. Cocomplete and
locally small; (co)limits pointwise from $\mathcal{S}$.

\medskip
\textbf{Yoneda embedding.}\quad $y : \mathcal{C} \hookrightarrow
\mathcal{P}(\mathcal{C})$, $y(x) = \mathrm{Map}_\mathcal{C}(-, x)$.

\medskip
\textbf{$\infty$-Yoneda lemma.}\quad
\begin{equation*}
  \mathrm{Map}_{\mathcal{P}(\mathcal{C})}(y(x), F) \;\simeq\; F(x),
  \qquad \text{natural in } x \in \mathcal{C}, F \in \mathcal{P}(\mathcal{C}).
\end{equation*}

\medskip
\textbf{Full faithfulness.}\quad $y$ induces weak equivalences
$\mathrm{Map}_\mathcal{C}(x, y) \xrightarrow{\sim} \mathrm{Map}_{\mathcal{P}(\mathcal{C})}(y(x), y(y))$.

\medskip
\textbf{Free cocompletion.}\quad For cocomplete $\mathcal{D}$:
\begin{equation*}
  \mathrm{Fun}^{\mathrm{cc}}(\mathcal{P}(\mathcal{C}), \mathcal{D}) \;\xrightarrow{\;y^*\;}\; \mathrm{Fun}(\mathcal{C}, \mathcal{D})
\end{equation*}
is an equivalence. The unique extension of $F : \mathcal{C} \to \mathcal{D}$
is $\bar F(P) = \mathrm{colim}_{(x, \xi) \in \mathcal{C}_{/P}} F(x)$.

\medskip
\textbf{Density.}\quad Every $P \in \mathcal{P}(\mathcal{C})$ satisfies
\begin{equation*}
  P \;\simeq\; \mathrm{colim}_{(x, \xi) \in \mathcal{C}_{/P}} \, y(x).
\end{equation*}

\medskip
\textbf{Bridge.}\quad With $\infty$-Yoneda available, the higher-CT
literature (Lurie HTT, Riehl–Verity, Cisinski, Kerodon, Land) is
accessible. Phase~2 of Volume~IV closes here.
\end{formalrestatement}
